\uuid{aRu0}
\titre{Racine}
\auteur{bodin ; hanani ; mzari}
\organisation{exo7}
\licence{CC BY-NC-SA 4.0}
\datecreate{2020-01-20}
\niveau{L1}
\module{Algèbre}
\chapitre{Polynôme, fraction rationnelle}
\sousChapitre{Racine, décomposition en facteurs irréductibles}
\difficulte{3}
\theme{Racines, factorisation, racine}
\exo7qcmid{lille-1.polynomes.024}

\contenu{
\begin{question}
\qclassification{105.03}
\qtags{motcle=racine}
Soit $P \in \Kk[X]$ un polynôme de degré $n\ge1$.
Quelles sont les affirmations vraies ?
\begin{answers}
    \good{$a$ racine de $P$ $\iff$ $X-a$ divise $P$.}

    \good{$a$ racine de $P$ de multiplicité $\ge k$ $\iff$ $(X-a)^k$ divise $P$.}

    \bad{$a$ racine de $P$ de multiplicité $\ge k$ $\iff$
    $P(a) = 0$, $P'(a)=0$, ..., $P^{(k)}(a)=0$.}

    \good{La somme des multiplicités des racines est $\le n$.}

\end{answers}
\begin{explanations}
$a$ racine de $P$ de multiplicité $\ge k$ $\iff$ $(X-a)^k$ divise $P$ $\iff$ $P(a) = 0$, $P'(a)=0$, ..., $P^{(k-1)}(a)=0$.
\end{explanations}
\end{question}
}
